% This is text associated with the open textbook "Open Electromagnetics"
% See immediately below for author, date
% This work is licensed under the Creative Commons Attribution-ShareAlike 4.0 International License. (CC BY SA 4.0) 
% To view a copy of the license, visit http://creativecommons.org/licenses/by-sa/4.0/

%ORB TITLE Decibel Scale for Power Ratio
%ORB AUTHOR 0        # S.W. Ellingson, Virginia Tech (ellingson@vt.edu)
%ORB DATE 20191115   # yyyymmdd
%ORB ABSTRACT About the Author (Ellingson)

\label{m0154_Decibel_Scale_for_Power_Ratio} % <-- Must match name of this file and module directory name.  Don't mess with this!
\index{decibel|(}
\index{dB|see{decibel}}

In many disciplines within electrical engineering, it is common to evaluate the ratios of powers and power densities that differ by many orders of magnitude.
These ratios could be expressed in scientific notation, but it is more common to use the logarithmic \emph{decibel} (dB) scale in such applications.

In the conventional (linear) scale, the ratio of power $P_1$ to power $P_0$ is simply
%
\begin{equation} 
G = \frac{P_1}{P_0} ~~~\mbox{(linear units)} 
\end{equation}
%
Here, ``$G$'' might be interpreted as ``power gain.''
\index{gain!power}
Note that $G<1$ if $P_1<P_0$ and $G>1$ if $P_1>P_0$.

In the decibel scale, the ratio of power $P_1$ to power $P_0$ is
%
\begin{equation}
\boxed{
G \triangleq 10\log_{10}\frac{P_1}{P_0} ~~~\mbox{(dB)}
}
\label{m0154_eGdef}
\end{equation} 
%
where ``dB'' denotes a unitless quantity which is expressed in the decibel scale.
Note that $G<0$~dB (i.e., is ``negative in dB'') if $P_1<P_0$ and $G>0$~dB if $P_1>P_0$.
%
%%%%%%%%%%%%%%%%%%%%%%%%%%%%%%%%%%%%%%%%%%%%%%%
%ORB HIGHLIGHT
\begin{mdframed}[backgroundcolor=yellow!20,nobreak=true] 
The power gain $P_1/P_0$ in dB is given by Equation~\ref{m0154_eGdef}.
\end{mdframed}
%%%%%%%%%%%%%%%%%%%%%%%%%%%%%%%%%%%%%%%%%%%%%%%

Alternatively, one might choose to interpret a power ratio as a loss $L$ with $L\triangleq 1/G$ in linear units, which is $L=-G$ when expressed in dB.
Most often, but not always, engineers interpret a power ratio as ``gain'' if the output power is expected to be greater than input power (e.g., as expected for an amplifier) and 
as ``loss'' 
if output power is expected to be less than input power (e.g., as expected for a lossy transmission line). 
%
%%%%%%%%%%%%%%%%%%%%%%%%%%%%%%%%%%%%%%%%%%%%%%%
%ORB HIGHLIGHT
\begin{mdframed}[backgroundcolor=yellow!20,nobreak=true] 
Power loss $L$ is the reciprocal of power gain $G$.  Therefore, $L=-G$ when these quantities are expressed in dB.
\end{mdframed}
%%%%%%%%%%%%%%%%%%%%%%%%%%%%%%%%%%%%%%%%%%%%%%%

%%%%%%%%%%%%%%%%%%%%%%%%%%%%%%%%%%%%%%%%%%%%%%%
%ORB EXAMPLE
~\\ \begin{example} 
\begin{mdframed}[backgroundcolor=brown!20,linecolor=brown!20]\raggedright
Power loss from a long cable.
\\~\\
A 2~W signal is injected into a long cable.  
The power arriving at the other end of the cable is 10~$\mu$W.  
What is the power loss in dB? 

In linear units:
%
\begin{equation}
G = \frac{10~\mbox{$\mu$W}}{2~\mbox{W}} = 5 \times 10^{-6} ~~~\mbox{(linear units)} \nonumber
\end{equation}
%
In dB:
%
\begin{flalign}
G &= 10\log_{10}\left(5 \times 10^{-6}\right) \cong -53.0~\mbox{dB} \nonumber \\
L &= -G \cong \underline{+53.0~\mbox{dB}} \nonumber
\end{flalign}
%
\end{mdframed}
% note figures will need to go between "\end{mdframed}" and "\end{example}"
\end{example}
%%%%%%%%%%%%%%%%%%%%%%%%%%%%%%%%%%%%%%%%%%%%%%%

The decibel scale is used in precisely the same way to relate ratios of spatial power densities for waves.  
For example, the loss incurred when the spatial power density is reduced from $S_0$ (SI base units of W/m$^2$) to $S_1$ is
%
\begin{equation}
L = 10\log_{10}\frac{S_0}{S_1} ~~~\mbox{(dB)}
\end{equation}
%
This works because the common units of m$^{-2}$ in the numerator and denominator cancel, leaving a power ratio.

A common point of confusion is the proper use of the decibel scale to represent voltage or current ratios.  
To avoid confusion, simply refer to the definition expressed in Equation~\ref{m0154_eGdef}.
For example, let's say 
$P_1 = V_1^2/R_1$ where $V_1$ is potential and $R_1$ is the impedance across which $V_1$ is defined.
Similarly, let us define 
$P_0 = V_0^2/R_0$ where $V_0$ is potential and $R_0$ is the impedance across which $V_0$ is defined.
Applying Equation~\ref{m0154_eGdef}:
%
\begin{flalign}
G &\triangleq 10\log_{10}\frac{P_1}{P_0} ~~~\mbox{(dB)} \nonumber \\
  &= 10\log_{10}\frac{V_1^2/R_1}{V_0^2/R_0} ~~~\mbox{(dB)}
\end{flalign}
%
Now, if $R_1=R_0$, then
%
\begin{flalign}
G &= 10\log_{10}\frac{V_1^2}{V_0^2} ~~~\mbox{(dB)} \nonumber \\
  &= 10\log_{10}\left(\frac{V_1}{V_0}\right)^2 ~~~\mbox{(dB)} \nonumber \\
  &= 20\log_{10}\frac{V_1}{V_0} ~~~\mbox{(dB)}
\end{flalign}
%
However, note that this is \emph{not} true if $R_1\neq R_0$.
%
%%%%%%%%%%%%%%%%%%%%%%%%%%%%%%%%%%%%%%%%%%%%%%%
%ORB HIGHLIGHT
\begin{mdframed}[backgroundcolor=yellow!20,nobreak=true] 
A power ratio in dB is equal to $20\log_{10}$ of the voltage ratio only if the associated impedances are equal.
\end{mdframed}
%%%%%%%%%%%%%%%%%%%%%%%%%%%%%%%%%%%%%%%%%%%%%%%

Adding to the potential for confusion on this point is the concept of \emph{voltage gain} $G_v$:
\index{gain!voltage}
%
\begin{equation}
G_v \triangleq 20\log_{10}\frac{V_1}{V_0} ~~~\mbox{(dB)}
\end{equation}
%
which applies regardless of the associated impedances. 
Note that $G_v=G$ only if the associated impedances are equal, and that these ratios are different otherwise. 
Be careful!

The decibel scale simplifies common calculations.
Here's an example.
Let's say a signal having power $P_0$ is injected into a transmission line having loss $L$.
Then the output power $P_1=P_0/L$ in linear units.  However, in dB, we find:
%
\begin{flalign}
10\log_{10}{P_1} &= 10\log_{10}{\frac{P_0}{L}} \nonumber \\
               &= 10\log_{10}{P_0} - 10\log_{10}{L} \nonumber 
\end{flalign}
%
Division has been transformed into subtraction; i.e., 
%
\begin{equation}
P_1 = P_0 - L ~~~\mbox{(dB)}
\label{m0154_edbd}
\end{equation}
%
This form facilitates easier calculation and visualization, and so is typically preferred.

Finally, note that the units of $P_1$ and $P_0$ in Equation~\ref{m0154_edbd} are not dB \emph{per se}, but rather dB with respect to the original power units.
For example, if $P_1$ is in mW, then taking $10\log_{10}$ of this quantity results in a quantity having units of dB relative to 1 mW.
A power expressed in dB relative to 1~mW is said to have units of ``dBm.''
\index{dBm}
For example, ``0~dBm'' means 0~dB relative to 1~mW, which is simply 1~mW.
Similarly $+10$~dBm is 10~mW, $-10$~dBm is 0.1~mW, and so on.

%%%%%%%%%%%%%%%%%%%%%

{\bf Additional Reading:}
\begin{itemize}
\item ``\href{https://en.wikipedia.org/wiki/Decibel}{Decibel}'' on Wikipedia.
\item ``\href{https://en.wikipedia.org/wiki/Scientific_notation}{Scientific notation}'' on Wikipedia.
\end{itemize}

\index{decibel|)}


